\begin{afterword}
\summaryhead

The post is about people who say that science ``doesn't really know anything'' because it admits it might be wrong. For them, the author says, a source is trusted wholly or not at all, so an admitted error ends all trust. The author names religion, where doubting scripture is a sin, as one obvious source of this habit, and suspects that school is another, since the teacher's statements must be believed and students' suggestions need not be.

In a public debate the author would try short lines: that science's strength is changing its mind, and that a belief one would give up for enough evidence cannot have probability one. In private the first lesson is that one can choose between better and worse without certainty. The claim that without certainty ``the moral relativists win'' fails, since certainty is sufficient for choosing but not necessary. Probability should mean calibration; people's ``absolute'' certainties are not all correct; and probability 1 would be infinite certainty. Scientists, the author suggests, should say ``We are not infinitely certain.''

\responsehead

The post's central point is correct. Choosing well needs knowledge of better and worse, not certainty about good and evil, and a belief held with probability one cannot be moved by any evidence. Statisticians know the second point as Cromwell's rule. The claim that people's ``absolutely certain'' statements are not all correct is also true; the author cites studies that show it in a book chapter, mentioned in the next day's post.

The post treats its opponent unfairly. The post opens with a real challenge: science changed its mind about gravity, so why trust it on evolution? That question contains an argument, not only a demand for certainty. Past theories were replaced, so present ones may be too. Philosophers call this the pessimistic induction and answer it with care. The post answers only the demand for certainty, and gravity never comes up again. A standard reply was in the Asimov essay quoted at the end of the previous day's post: good theories are refined, and are ``not so much wrong as incomplete.''

The post also explains the opponent's view by where it supposedly comes from. Its two sources, religion and school, are introduced with ``One obvious source'', ``I suspect'' and ``I fear'', and nothing in the post tests them. The opponents are described rather than quoted. They are called ``the unenlightened ones'', and a boast, ``I'm perfect'', is put in their mouths. A parenthesis treats people who doubt Bayesian priors the same way. Whether priors can be argued with is an open problem that divides Bayesians, and the post explains the doubt as a habit learned in class.

The calibration section states its conclusion more firmly than its evidence allows. The claim that commitment ``doesn't even behave monotonically'' rests on sentences hedged with ``If anything'' and ``probably''. The one data set the author had quoted elsewhere shows accuracy rising with stated confidence at the top of the scale. Stated odds are not the same as the emotional commitment the post describes, so the data do not refute the post, but they do not support it either.

In short: a correct point, that one can choose without certainty, made against an opponent who is described rather than quoted, whose question about a science that changed its mind is answered only as a demand for certainty, and whose view is traced to religion and school by guesses the post does not test.
\end{afterword}
